% ================================================================
%  第 6 章 梦与做梦者 后半 (原书 p.147-152, 6.9-6.11)
%  源文: _work/src1/task-ch06b.md
%  注意:首段紧接 part1-ch06a.tex 末句 “唉!竟然不得不”, 两者之间不可插入空行
% ================================================================
% ---------------- p.147 ----------------
学到(或者:让人领教)这种事——谁会相信呢——真得是个混账(我差点要说:废物,对不起!),才会那样对待活生生的人!说到底,这跟另一副调子也差不了多少:只消把“废物”换成“混账”,把“让人那样对待”换成“那样对待”,戏法就变完了!而荣誉嘛,当然,毫发无损——对善事业的旗手来说!

\origpage{29} 从这里清楚显出的,是我与种种轻蔑态度的默契。它至少上溯到五十年代初,也就是我在嘉当和他朋友们那里受到善意接待之后的那些年份。后来,当轻蔑在四面八方变得司空见惯时,我之所以“什么也看不见”,是因为我不想看见——并不比那桩孤立而格外扎眼的事里更想看:那一回,真得下足功夫,才能装作既看不见、也觉不出!

这份默契与我新的身份紧密共生——一个受尊敬的群体成员的身份,有功者、数学强者的群体。我记得我特别满意,甚至骄傲:在我为自己选定、又把我增选进去的那个世界里,起作用的不是社会地位,甚至(才不是呢!)也不单是名声——名声还得是挣来的;你尽可以是大学教授、院士或随便什么,可你要只是个平庸的数学家(可怜的家伙!),那你就什么也不是;唯一算数的是功绩,是深刻而独到的想法,是技术上的娴熟,是广阔的视野,诸如此类!

这套功绩主义,我曾毫无保留地认同于它(而它当然始终隐而不宣、不曾明言),可它在我身上到底还是挨了狠狠一击——就在我说过的1970年那场著名觉醒之后的日子里。而且我也不能肯定,它从那一刻起就消失了、没留下痕迹。要做到那一点,大概得靠我清清楚楚地在自身中把它识别出来;可我当时似乎主要是在别人身上谴责它。说起来,最早提醒我注意这套意识形态的,正是谢瓦莱,还有德尼·盖吉(我也是通过“生存”认识他的)——他们管它叫“功绩主义”(méritocratie),或者类似的名字——以及它里面那种暴力、那种轻蔑。谢瓦莱对我说(那大概是我们头一次在他家里见面、谈“生存”的时候),正是因为这个,他才再也受不了布尔巴基的气氛,索性不再踏进去。回想起来,我确信他当时一定早已看出:我实实在在是这套意识形态的一分子,也许甚至看出它在我的某些角落里还留着痕迹。但我不记得他曾有过任何暗示。也许这一回,他仍然宁愿让我自己去把那些他替我勾出的 i 点上点——而我一直等到今天才把它们点上。晚做总比不做强!

\section{(13) 力与厚度}

\origpage{30} 我讲的那桩事,很可能也标志着我内心一次倒向的时刻:转向对功绩兄弟会或多或少无条件的认同,而以那些被视为废物的人为代价——或者,按早几代人的说法,只消说“没有天分”(这个词到我那时候已经不流行了):那些平庸、暗淡的人——顶多算是“共鸣箱”(韦伊在什么地方这样写过),供真正要紧的人的大想法在其中回响……我的记忆常常像个掘墓人,连当时耗掉大量心理能量的片断也一并埋掉;可它偏偏留下了这一段,而这一段跟任何别的记忆都没有直接关联,面目又那么平淡无奇——单凭这一点,那个“倒向”当时确曾发生的感受,就变得可信了。

不到五年前的一次沉思中,我终于意识到:那种“我们,伟大而高贵的精神……”的意识形态,以一种格外极端、格外凶猛的形式,从我母亲童年起就在她身上肆虐,并且主宰着她与他人的关系——她乐于从自己的高处俯视那些人,带着一种往往是鄙夷的、甚至轻蔑的怜悯。
% ---------------- p.148 ----------------
\origpage{31} 再说,我对我的父母怀着毫无保留的敬佩。在那个著名的“数学共同体”之前,我认同的第一个、也是唯一一个群体,就是缩小到母亲、父亲和我三个人的家庭——而我有幸被母亲认可为配得上他们做我的父母。这就是说,轻蔑的种子想必从我的童年起就已播进我这个人里了。也许时机已经成熟,可以追一追这些种子在我的童年和成年生活中经历的种种变故,以及其中一些种子长成的收成——幻觉的收成、孤立的收成、冲突的收成。但这里不是做这件事的地方,此处我追着一个更有限的目标。我想我可以说:这种轻蔑态度在我一生中,从未达到我在母亲一生中见到的那种激烈、那种毁灭性的力量(那是我肯费心去看我父母一生的时候看到的——在母亲去世二十二年之后,在父亲去世三十七年之后)。但现在正是时候,过了这村就没这店:在这里至少仔细考察一下,这种态度在我的数学生涯中占了什么位置。

在这之前,为了把上一段讲到的那桩事放进它的一般背景里,我想强调这样一个事实:在我的五十年代记忆里——甚至更晚的记忆里——它完全是孤立的。即便在今天,虽然我发觉在我曾经所属的那个圈子里,某些最基本的礼貌和对他人的尊重,有时正以一种叫人错愕的方式被侵蚀\fnoteP{8}{(10)}{比方说,我已数不清有多少封信——谈的既有数学问题,也有实际问题或私事——寄给我当作朋友的同事或旧学生,却从未收到回音。看来这并非只是专门留给我一个人的特殊待遇,而确实是一种风气变化的迹象——据一些同方向的反响来看。(诚然,那些反响涉及的情形是:寄数学信件的人,收信的那位知名数学家并不认识他……)} (10),老板对学生那种直接而不加掩饰的轻蔑,也还应当是相当罕见的事。至于五十年代,我很少有记忆指向那种情形:一个有名望的人物周围笼罩着恐惧,或者有谁摆出轻蔑、哪怕只是鄙夷的态度。要往这个方向翻找,我可以说:头一回被请到南锡迪厄多内家做客时——他对我一向有着那份体贴入微的和气——这个文雅亲切的人谈起他的学生来,那副样子真叫我有点发懵:简直全是些蠢货!给他们上课是件苦差,而他们显然什么也听不懂……1970年以后,我听到从阶梯教室那一侧传来的回声,才知道迪厄多内确实让学生们害怕。可是,尽管他以见解果决、并且常常以雷鸣般的直率把它们讲出来而闻名,我却从未见他有过伤人、羞辱人的举动——即使当着他不怎么看重的同事的面,即使在他那些传奇式的雷霆大怒之际;而那怒气,来得快,消得也一样快、一样轻易。

我不认同迪厄多内关于他学生的那些说法,可也没有同他的态度拉开距离——那态度被呈现为世上最理所当然的事,对一个热爱数学的人来说几乎是不言自明的。有这位年长同行满怀善意的权威作后盾,那种态度当时在我眼里,至少是人们在面对学生和教学任务时可以有理由采取的态度之一。

看来,迪厄多内和我一样,都浸透了这套功绩主义;但当我们面对一个血肉之躯时,这种主义的隔离效果在很大程度上就被抵消了:单是那个人的在场,就默默地向我们提醒一些比所谓“功绩”更本质的现实,并重新接上一根被遗忘的纽带。我们的大多数同事或朋友想必也一样——他们浸染那种如此普遍的优越感,并不比迪厄多内或我少。今天,他们中的许多人肯定还是这样。

% ---------------- p.149 ----------------
\origpage{32} 韦伊同样有让学生害怕的名声;而在五十年代我的微观世界里,他是唯一一个让我觉得连那些地位(或者干脆说性情)比较谦卑的同事也怕他的人。他时不时会摆出那种不容回嘴的居高临下,足以叫最牢靠的自信也发慌。我的敏感也在其中起了作用,于是有那么一两次,我们闹了阵短暂的别扭。我在他的做派里没有察觉出一丝轻蔑的意味,也没有察觉出存心伤人、压人的意图;那更像是被惯坏的孩子的作派,以叫人难堪为乐(有时还带点促狭),仿佛以此使自己确信他握有某种权力。再说,他对布尔巴基小组有一种真正惊人的支配力;有时他给我的印象是,他拿指挥棒指挥着它,有点像幼儿园女教师指挥一群乖孩子。

在五十年代,我只记得另外一次,感到轻蔑以一种粗暴、不加掩饰的方式表达出来。它来自一位外国同事兼朋友,年纪跟我差不多。他的数学力量非同寻常。几年前,这股力量已经相当明显了,而我当时却被他对那位大教授的顺从(在我看来近乎谄媚)所触动——那时他还只是那人的一名谦卑助手。他出众的才能很快为他赢得了国际声誉,以及一所格外显赫的大学里的一个关键职位。他在那里统治着一小支助手兼学生的队伍,那副绝对的派头,显然跟他当年的老板统治他和他的同伴们时一样。我问他(如果我记得不错)有没有几个学生(言下之意:跟着他做出好东西的学生),他带着一副假作随意的神气(我把它译成法语)答道:“十二件!”——所以“件”就是他用来指他的学生和助手的词。一个数学家同时有这么多学生在他指导下做研究,确实少见——而我的这位对话者想必暗自引以为傲,又想用那副漫不经心的神气把它藏起来,仿佛在说:“哦,不过十二件,连提都不值一提!”那大概是1959年前后。我当时想必已经有了一层不错的甲壳,可我还是感到一阵恶心!我准是当场就用某种方式对他说了,而且我不认为他为此记恨我。也许他与学生的关系并不像他那说法所能让人设想的那样阴森(我没有得到过他任何一位学生的证词),而他只是落进了自己那孩子气的欲望的陷阱:想在我面前炫耀他全部的荣光。回头看,我看得出这桩事必定是我们关系的一个转折点,而先前的那些关系是友谊的关系——我在他身上感到某种脆弱,还有一种细腻,它们吸引着我亲切的同情。这些品质被磨损了、被腐蚀了,被他那位重要人物的地位——受人敬佩、叫人害怕的地位。这桩事之后,我心里对他留下了一种不适——说到底,我觉得自己不属于他那个世界……

\origpage{33} 然而我们确实属于同一个世界——而且,我并不比他更觉察到:我也在变厚。关于这一点,我留下了一段鲜明的记忆,那是在1958年的爱丁堡国际数学家大会上。从上一年的黎曼–罗赫定理工作起,我就被晋升为大明星了;而且(当时我并不曾用明确的话对自己说过)我也是大会的明星之一。(我在会上作了一个报告,讲同年概形理论的强劲起步。)希尔策布鲁赫(当日的另一位明星,带着他自己的黎曼–罗赫定理)致开幕辞,为的是向那年即将退休的霍奇致敬。说到某处,希尔策布鲁赫暗示:数学主要靠年轻人的工作,而不是靠年长数学家的。这在年轻人占多数的大会厅里激起了一片轰然的赞同。我满心欢喜,当然也完全同意:我刚满三十岁,还满可以算年轻,而世界是属于我的!一兴奋,我准是大声叫喊,还在桌子上猛捶了好几下。碰巧我坐在霍奇夫人(Lady Hodge)旁边——她就是那位
% ---------------- p.150 ----------------
\origpage{34} 本该在这次场合受到致敬的卓越数学家的妻子,而他今年就要退休了。她睁大眼睛转向我,对我说了几句话,我已不记得了——但我准是从她惊讶的眼睛里看到了映出的那份没有分寸的厚度:它刚刚就在这位已近暮年的女士面前毫无节制地发作了一通。那时我感到了一点什么,而“羞耻”这个词给它的形象也许是走了样的——那更像是一条朴素的真相,关于当时的我这个人。那天,我大概再没有在桌子上猛捶过。

\section{(14) 恐惧的诞生}

我猜想,正是在那个时候——当我(并非刻意去求)开始被看作数学界的一个明星的时候——某种恐惧也必定开始围绕着我这个人,对许多我不认识或不太认识的同事来说。我是猜的:我无法用一个确切的记忆把它定位,无法用一幅曾击中我、并留在我记忆里的画面来定位,就像前面讲的那桩事那样(那想必标志着我头一回在接纳我的圈子里遇上轻蔑)。

这事想必是不知不觉地发生的,没有引起我的注意,也没有表现为某个特别而典型的、会被记忆保留下来的事件;而记忆打在它上面的那道光,也许跟打在那另一桩事上的一样,是刻意平淡的。我对那些过渡年月的记忆“整块地”还给我的,是这样一件事:前来接近我的人——无论是在我的讨论班之后,还是在布尔巴基讨论班之类的聚会上、在某个学术会议或大会上——往往得克服一种怯场,而在我们交谈(如果有交谈的话)的过程中,这种怯场多多少少还看得出来。当交谈超过几分钟,这份局促往往就在我们谈话、谈话热起来的过程中渐渐消失。有时——很少——也会发生这种情形:局促一直保持着,以至于成了交流的真正障碍,甚至在最不涉及个人的层面、在一场数学讨论的层面上;而我当时就模糊地感到对面有一种无力的痛苦,一种对自己恼火的痛苦。我说这一切时,并不是真的“记得”,而像隔着一层雾;可这层雾仍然把一些想必被记录过、又多半随即被排遣掉的印象还给了我。要我说出这份局促——这种恐惧的表现——何时出现,除了靠猜,我实在办不到。

我不认为这种恐惧源自我这个人,也不认为它只限于某种态度、某些使我区别于同事的举止。若真是那样,我想我在七十年代初就该听到些回声了——那时我走出一个自己一直扮演的角色,正是明星、“大老板”的角色。我相信,受人畏惧的是这个角色,而不是我这个人。而这个角色,连同它那圈与尊重毫无共同之处的恐惧光晕,在五十年代初并不存在,还没有存在——至少在我1948年一遇上就接纳了我的那个数学圈子里并不存在。

再说,在1970年那场“觉醒”之前,我根本不会想到把这种怯场、这种我有时在不那么亲近的同事身上遇到的局促叫作“恐惧”。它一显露,我自己就感到不自在,于是尽力去消除它。有一件事很值得注意,也典型地说明我亲爱的微观世界对这类事多么不关心:在我属于这个圈子的二十年里,我不记得有过一次,这个问题在我和某位同事之间被谈起,或者被别的人当着我的面谈起!\fnoteP{9}{(11)}{\emph{阿尔多·安德烈奥蒂(Aldo Andreotti)、约内尔·布库尔(Ionel Bucur)}\par{}当然,我这边也未必没有遗漏——何况我那时候那副十足“书呆子”的脾气,大概也不会鼓励别人跟我谈这类事,更不会促使我去记住一场很可能确实有过的、朝这个方向的谈话。可以肯定的是,至少可以说,这个问题被谈起的场合想必极其罕见(甚至不用“恐惧”这个名字去称呼它……),今天想必也一如其旧,尤其是在“上流社会”里。\par 在那个世界里的众多朋友中,除了谢瓦莱(他至少在六十年代期间必定意识到了这种恐惧的气氛),我觉得确实清楚察觉过它的,只有另外一个人,就是阿尔多·安德烈奥蒂。我认识他,以及他的妻子芭芭拉(Barbara)和他们的一对双胞胎孩子(当时还很小),是在1955年(我想是在芝加哥韦伊家的一次晚会上)。我们一直关系很密,直到1970年“大转折”的时候——那时我离开了曾属于我们的那个圈子,同他们也稍稍失去了联系。阿尔多极其敏感,而这份敏感在与数学、与像我这样的“书呆子”的往来中丝毫没有磨钝。他身上有一种对接近他的人自然生发的同情之才。这使他跟我所认识的数学圈内、乃至圈外的所有别的朋友都不一样。在他那里,友谊总是压过共同的数学兴趣(那些兴趣并不缺),而他也是极少数与我稍稍谈过我的一生、他也谈过他自己的一生的数学家之一。他的父亲和我的父亲一样是犹太人,他也因此在墨索里尼的意大利吃过苦,正如我在希特勒的德国吃过苦。我见他总是乐于鼓励和支持年轻研究者——那是在一个越来越难以被当权圈子接纳的气候里。他那种自然的关切,总先落到人身上,而不是落在什么数学“潜力”或名声上。他是我有幸遇到的最叫人亲近的人之一。\par 说到阿尔多,又勾起了对约内尔·布库尔的记忆——他也是出人意料地、未到年纪就被带走了;而且和阿尔多一样,我想,人们更多是把他当作一位乐于重聚的朋友来怀念,而不是当作数学讨论的伙伴。人能从他身上感到一种善良,旁边还有一份不多见的谦逊,一种不断退到后面的倾向。这样一个从不看重自己、也从不给人压力的人,最后竟当上了布加勒斯特大学理学院院长,真是个谜;也许是因为他压根儿没想过要推辞那些他远不渴求、却由同事们或政治当局放到他肩上的职务——那肩膀,说起来,是结实的。他是农民的儿子(在一个“阶级标准”很重要的国家里,这想必起了作用),身上带着农民的通达与朴素。他必定察觉到围绕着名人的那种恐惧,但想必也觉得这事理所当然,是权力地位的自然属性。不过,我不认为他本人曾让任何人畏惧——当然也不曾让他的妻子弗洛丽卡(Florica)、他们的女儿亚历山德拉(Alexandra),或他的同事、学生畏惧;而我能听到的反响,也正是这个方向。} (11)那层替我充作记忆的“雾”,也没有还给我
% ---------------- p.151 ----------------
\origpage{35} 任何印象,像是这类情形曾在我身上引起过有意识或无意识的满足。我不认为在意识层面有过;但我也不敢断言,在最初那些年里,我在无意识层面不曾偶尔被它拂过。若有,那也必定是短暂的,不曾反映到行为里、起一种把局促固定下来的作用。当然,这并不是说我演的那个角色里没有掺进我的自负!但即使我在这角色上不计本钱地投入,当时驱动我这份自我的,并不是想给“普通同事”留下深刻印象的野心,而是不断超越自己,以逼出我的“同侪”不断更新的敬意——而且,也许首先是那些在我还无从施展之前就信任我、把我当作他们中一员的年长者。看来,我面对自己成为对象的恐惧时所持的内心态度——尽力不去理会它,同时在它显露之处勉力把它消除——这态度可以看作整个六十年代我所属那个圈子(“微观世界”)里的典型态度。

在此后过去的十年到十五年里,局面又大大恶化了——至少从那个世界时不时传到我这里的迹象看是如此,从我得以近距离目击、甚至有时参与其中的情形看也是如此。不止一次,在我旧日的朋友或学生中那些与我曾最亲近的人那里,我撞见了轻蔑那种熟悉而无可辩驳的迹象;撞见了那种(表面上“无缘无故”的)使人泄气、羞辱人、压垮人的意志。一阵轻蔑之风起来了——我说不清是什么时候——吹拂着那个曾于我亲切的世界。它吹着,不管什么“功绩”或“过”,用它的气息烧焦卑微的天职,也烧焦最美好的激情。我旧日的伙伴们,各自带着“自己人”,受着坚固城墙的保护,安稳地(像我从前那样)坐在那层裹着各自的、毡子般悄然无声的恐惧里——他们中可有一个人感到这阵风?我确切地认识一个,只有一个,在我的旧友当中:他感到了,还跟我谈起过,只是没有叫出它的名字。还有另外一位,有一天也像是身不由己地察觉到了它,随即赶忙在翌日\fnoteP{10}{(12)}{“翌日”一词在此应按字面理解,不是隐喻。} (12)就把它忘掉。因为感到这阵风并承担它,对我的一位旧友来说,
% ---------------- p.152 ----------------
\origpage{36} 正如对我自己一样,也就是接受把目光落到自己身上。

\section{(15) 收获与播种}

我不打算——我再也不会打算——为一阵吹着的风而愤慨了,既然我已看清:我与这阵风并非无关,尽管我身上的一份自负一心要让我相信相反的事。而即便我与它无关,我的愤慨对受辱的人和施辱的人也都只会是一份太可怜的献礼——而这两种人,我一样爱过。

我与这阵风并非无关:在我为自己选定的那个世界里,我与轻蔑、与恐惧有过默契。对这些污迹,就像对许多别的污迹一样,闭上眼睛对我正方便——在我的职业生活中如此,在我的家庭生活中也一样。在这两者里,我收获了我所播种的——也收获了别人在我之前或与我一同播种的,既包括我的父母(以及我父母的父母……),也包括我那些“旧日的新朋友”。而今天,收获这些发了芽的播种的,还有我之外的别人:既有我的孩子们(以及我孩子的孩子),也有我今天的某个学生——他被我旧日的某个学生轻蔑地对待。

谈到播种与收获,我心里既没有苦涩,也没有认命,更没有自怜。因为我已懂得:即便在苦涩的收成里,也有一种实体的肉,只消我们肯吃,就能拿它来养自己。当这份实体被吃下,成了我们血肉的一部分,苦涩就消失了——它原来只不过是我们面对一份为我们准备的食物时的抵抗的征兆。

我也知道:没有哪一次收获不同时是对另一些收获的播种——而那些收获往往比前一次更苦。当那条看来没有尽头的链条——漫不经心的播种与苦涩的收获,代代相传、代代重演——出现在眼前时,我身上仍有某处会揪紧。但我不再像面对一种残酷而无从逃避的宿命那样被压垮、起而反抗,更不再像从前那样做它心甘情愿而盲目的囚徒。因为我知道,在我遇到的一切里都有一种滋养的实体,无论播种的是我的手还是别人的手——只消我去吃,去看它变成认识。对我的孩子们,对所有我爱过的人、以及此刻我爱着的人,情形也一样:当他们收获我在自负而漫不经心的年月里播下的东西,或者收获我直到今天偶尔还在播下的东西时,也是如此。
